% concatenated from main.tex
%------------------------------------------------------------------------------
% Beginning of journal.tex
%------------------------------------------------------------------------------
%
% AMS-LaTeX version 2 sample file for journals, based on amsart.cls.
%
%        ***     DO NOT USE THIS FILE AS A STARTER.      ***
%        ***  USE THE JOURNAL-SPECIFIC *.TEMPLATE FILE.  ***
%
% Replace amsart by the documentclass for the target journal, e.g., tran-l.
%\documentclass[a4paper,11pt,reqno]{amsart}
%
\documentclass[a4paper,12pt,reqno]{amsart}
\usepackage{amsthm, mathrsfs,amssymb,amsmath,}
\usepackage{enumerate}
\usepackage{mathtools}
\usepackage{breqn}
\makeatletter
\@namedef{subjclassname@2010}{%
	\textup{2010} Mathematics Subject Classification}
\makeatother

\frenchspacing

\setlength{\textwidth}{15.2cm}
\setlength{\textheight}{22.7cm}
\setlength{\topmargin}{0mm}
\setlength{\oddsidemargin}{3mm}
\setlength{\evensidemargin}{3mm}
%\setlength{\headsep}{2.9cm}
\setlength{\footskip}{1cm}
\usepackage{hyperref}
%\usepackage[colorlinks,citecolor=blue,pagebackref,hypertexnames=false]{hyperref}
\allowdisplaybreaks
\numberwithin{equation}{section}

\usepackage{graphicx}
\usepackage{enumitem}
\renewcommand{\labelenumii}{\Roman{enumii}}
\def\theequation{\arabic{section}.\arabic{equation}}
\renewcommand{\arraystretch}{2.5}


%\documentclass{amsart}
%\usepackage{graphicx} 
\newtheorem{theorem}{Theorem}[section]
\newtheorem{lemma}[theorem]{Lemma}
\usepackage{xcolor}

\theoremstyle{definition}
\newtheorem{definition}[theorem]{Definition}
\newtheorem{example}[theorem]{Example}
\newtheorem{proposition}[theorem]{Proposition}
\newtheorem{corollary}[theorem]{Corollary}
\newtheorem{note}[theorem]{Note}
\newtheorem{xca}[theorem]{Exercise}

\theoremstyle{remark}
\newtheorem{remark}[theorem]{Remark}

\numberwithin{equation}{section}


\begin{document}
	
	\title[Set Valued Riemann-Liouville integral and some Regular Selections]{Set Valued Riemann-Liouville integral and some Regular Selections}
	
	%    Information for first author
	
	
	
	
	
		\author{Subhash Chandra}
	%   Address of record for the research reported here
	\address{Harish-Chandra Research Institute, A CI of Homi Bhabha National Institute, Prayagraj, India}
	%    Current address
	
	\email{sahusubhash77@gmail.com}
	\author{Syed Abbas}
\address{School of Mathematical and Statistical Sciences, Indian Institute of Technology Mandi\\ Kamand (H.P.)- 175005, India}
\email{sabbas.iitk@gmail.com, abbas@iitmandi.ac.in ( Email of corresponding author)}
	
	
	
	

	%\date{\today}
	\subjclass{Primary: 28B20, 26A33.}
	
	\keywords{Riemann Liouville integral, Set valued mappings, Bounded variation, Regular selections}
	
	\begin{abstract}
In this article, we introduce the notion of the Riemann-Liouville fractional integral of set-valued mappings via integrable selections. We establish fundamental properties of this fractional integral, including convexity, boundedness, and continuity with respect to the Hausdorff metric. The investigation of preservation of regularity under fractional integration with respect to the Hausdorff metric is given. We show that bounded variation and Lipschitz continuity of a set-valued mapping are inherited by its Riemann-Liouville fractional integral.  We discuss the existence of regular selections for the fractional integral under the corresponding regularity assumptions on the original mapping. In the scalar case, we further identify extremal selections given by the pointwise minimum and maximum of the fractional integral and show that they possess the same regularity properties. Finally, we discuss possible applications in differential inclusion and directions for future research.
    

	\end{abstract}
	\maketitle
	
	
	
	
	\section{Introduction}
   The Riemann-Liouville integral is the first step towards defining fractional derivatives and hence fractional differential equations. This integral is a natural generalization of the regular n-fold integral of a suitable function. It is a fundamental operator in fractional calculus, which allows integration of non-integer values of the integration order $\alpha >0.$ This operator was developed through several contributions from mathematicians, most notably Liouville and Riemann, with recent formalizations by Sonin and Letnikov \cite{pod}. One can see that while classical calculus is limited to local, integer-order operations, the Riemann-Liouville integral provides non-local and memory-dependent behaviors. This enables us to model systems where past states influence the present, such as hereditary or anomalous processes. There is much theoretical and applied work in the field of fractional differential equations; for more details, we refer to \cite{abbas2, jehad1, li2}

    The set-valued integrals and differential equations are very important in several areas of engineering and physical sciences. One natural question to ask is whether the Riemann-Liouville integral can be extended to a set-valued Riemann-Liouville integral. So, the set-valued Riemann-Liouville integrals play an important role in extending fractional calculus to tackle imprecisions, uncertainties, or multiple possibilities in mathematical modeling. In the set-valued context, it generalizes this to functions that map to sets (subsets of $  \mathbb{R}  $ or higher dimensions) rather than single points. Let $\mathcal{K}_c(\mathbb{R})$ compact convex subset, for a set-valued function $F: [a, b] \to \mathcal{K}_c(\mathbb{R})$, the Riemann-Liouville fractional integral of order $\alpha > 0$ is often defined selection-wise. It can also be defined via embeddings like the Aumann integral, but in interval-valued settings, it is computed component-wise. 

In real-world systems, measurement mistakes, unpredictability, or inadequate data may make it impossible to determine parameters with precision. These are represented as sets via set-valued (including interval-valued) techniques. Set-valued analysis is a branch of mathematics that deals with functions mapping points to sets rather than single values. This setup is important in several areas such as optimization, control theory, differential inclusions, and economics. In recent times, the application in the area of reinforcement learning, where uncertainty or robustness leads to set-valued policies or dynamics has been explored for details, we refer to \cite{ergen}. Set-valued integrals, which apply classical integration to multifunctions, and selections, which select single-valued functions from a set-valued map while maintaining continuity or measurability properties, are key ideas. Such setup enables fractional integrals to investigate dynamics in fields like rheology (material deformation), viscoelasticity (memory effects in materials), diffusion processes, control theory, and mathematical economics. When inputs are intervals rather than points, they allow for robust optimization. These integrals are key for solving Cauchy-type problems in fractional differential equations with set-valued right-hand sides, which is called differential inclusion. We provide the application of set-valued selection in differential inclusion. As it is more general and applicable, it comes with a cost in terms of additional restriction. The set-valued Riemann-Liouville integrals present several theoretical and practical difficulties due to operational complexities, existence-uniqueness issues, and integrability constraints. These challenges actually highlight the need for more research with relaxed assumptions and exploring broader set-valued extensions beyond intervals.
    
\par

Historically, set-valued integration emerged in the 1960s with works by Aumann and Debreu, building on earlier vector measure theory. This survey article gives a nice overview and important papers on this topic, focusing more on mathematical developments and especially connections to RL where relevant; for details, we refer to \cite{anca}. For set-valued selection, we refer to \cite{hess}. The foundational paper defining the Aumann integral via selections is given in \cite{aumann1965integrals}. The idea of  measures and integrals for multifunctions is discussed in \cite{art}. Readers interested more in set-valued integrals, selections, and viability theory may see \cite{aubin1}. Lupulescu in \cite{lup} introduced the first systematic framework for fractional calculus on interval-valued functions. Several important properties, such as monotonicity, linearity, and semigroup composition, have been established. Abbas et al. \cite{azam} discussed fractional differential inclusion with multipoint boundary condition. Fixed-point techniques are utilized to study set-valued inclusions. 
Motivated by these developments, it is natural to ask whether the Riemann-Liouville fractional integral can be meaningfully extended to set-valued mappings. The next section covers the basic terminologies required to establish our motive. Section 3 is dedicated to our main findings and application. In Section 4, we give some important future directions related to the dimension theory of a set valued Riemann-Liouville fractional integral.

\pagebreak

\section{Preliminaries} \label{Prelim}

In this section, we lay out the background material related to this article.
\begin{definition} 
   Let $a<b$ and let $f\in L^{1}([a,b])$. For $\rho>0$, the
\emph{Riemann-Liouville fractional integral} of order $\rho$ of $f$ is
defined by
$$
{}_{a}J^{\rho}f(u) = \frac{1}{\Gamma(\rho)}
\int_{a}^{u}(u-t)^{\rho-1}f(t)\,dt,
\qquad u\in[a,b], $$
where $\Gamma(\cdot)$ denotes the Euler gamma function.
\end{definition}

\begin{definition}
Let $F:[a,b]\rightrightarrows \mathbb{R}$ be a set-valued mapping, then the set
\[Gr(F)=\left\{(u,w): u\in[a,b] \text{ and } w\in F(u)\right\},\]
is characterized as a graph of $F$. 
\end{definition}
\begin{definition}
Let $\mathcal{K}(\mathbb{R})$ denote the family of all nonempty compact 
subsets of $\mathbb{R}$. For $A,B\in\mathcal{K}(\mathbb{R})$, the 
\emph{Hausdorff distance} between $A$ and $B$ is defined by
\[
H_d(A,B)
:=
\max\left\{
\sup_{a\in A}\inf_{b\in B} |a-b|,\;
\sup_{b\in B}\inf_{a\in A} |a-b|
\right\}.
\]
Since $A$ and $B$ are compact subsets of $\mathbb{R}$, $H_d$ is finite and 
defines a metric on $\mathcal{K}(\mathbb{R})$.
\end{definition}

\begin{definition}
    Let $F:[a,b]\rightrightarrows\mathbb{R}$ be a set-valued mapping with nonempty
compact values. We say that $F$ is \emph{continuous} at $u_0\in[a,b]$ 
(with respect to the Hausdorff distance $H_d$) if 
\[
\lim_{u\to u_0} H_d\ \left(F(u),F(u_0)\right)=0.
\]
Equivalently, for every $\varepsilon>0$ there exists $\delta>0$ such that
\[
|u-u_0|<\delta
\quad \Longrightarrow \quad
H_d\!\left(F(u),F(u_0)\right)<\varepsilon .
\]
If this holds for every $u_0\in[a,b]$, then $F$ is said to be continuous on $[a,b]$ with respect to $H_d$.
\end{definition}
\begin{definition}
Let $F:[a,b]\rightrightarrows\mathbb{R}$ be a set-valued mapping with 
nonempty compact values. For a partition 
$\mathcal{P}=\{a=u_0<u_1<\dots<u_n=b\}$ of $[a,b]$, define
\[
V(F,\mathcal{P})
:=
\sum_{i=1}^{n} H_d\big(F(u_i),F(u_{i-1})\big).
\]
The \emph{total variation} of $F$ on $[a,b]$ is
\[
V(F,[a,b])
:=
\sup_{\mathcal{P}} V(F,\mathcal{P}),
\]
where the supremum is taken over all partitions $\mathcal{P}$ of $[a,b]$.
We say that $F$ is of \emph{bounded variation} on $[a,b]$ 
(w.r.t.\ $H_d$) if $V(F,[a,b])<\infty$.
\end{definition}
\begin{definition}
Let $F:[a,b]\rightrightarrows\mathbb{R}$ be a set-valued mapping with 
nonempty compact values. We say that $F$ is \emph{Lipschitz on} $[a,b]$ 
(with respect to $H_d$) if there exists a constant $L\ge 0$ such that
\[
H_d\big(F(u),F(v)\big)
\le L\,|u-v|
\qquad\text{for all }u,v\in[a,b].
\]
Any such $L$ is called a \emph{Lipschitz constant} of $F$ (with respect 
to $H_d$).
\end{definition}
\begin{definition}
A set-valued mapping, $F:[a,b]\rightrightarrows \mathbb{R}$ is said to be non-negative if for every $u\in [a,b]$ and $w\in F(u)$, we have $w\geq 0.$
\end{definition}



\begin{definition}\label{Bound}
A set-valued map, $F:[a,b]\rightrightarrows \mathbb{R}$ is said to be bounded, if it satisfies
\[\sup_{u\in [a,b]}H_{d}(F(u),\{0\})<\infty.\]
\end{definition}
\begin{note}\label{nt2.8} \label{NOTE}
Since, 
\begin{align*}
   \sup_{u\in [a,b]}H_d(F(u),\{0\})=&\sup_{u\in [a,b]}\left\{ \max \left\{\sup_{x\in F(u)}\inf_{y\in {0}}\lvert x-y \rvert,~ \sup_{y\in {0}}\inf_{x\in F(u)}\lvert y-x \rvert\right\} \right\}\\
   =& \sup_{u\in [a,b]}\left\{ \max \left\{\sup_{x\in F(u)}\lvert x \rvert,~ \inf_{x\in F(u)}\lvert x \rvert\right\} \right\}\\
   =& \sup_{u\in [a,b]} \sup_{x\in F(u)} \lvert x \rvert.
\end{align*}
Therefore, in view of Definition \ref{Bound} a function $F$ is bounded if it satisfies,
\begin{equation}\label{Eq1}
  \sup_{u\in [a,b]}~\sup_{x\in F(u)} \lvert x \rvert < \infty.  
\end{equation}
\end{note}

\begin{remark}
In view of Note \ref{nt2.8}, one can observe that unlike the single-valued map, a constant set-valued map need-not be bounded. 

For example, take $F:[0,1]\rightrightarrows \mathbb{R}$ is a set-valued constant map which is defined as $F(u)=[0,\infty)$ for each $u\in [0,1]$, then $F$ is not bounded as $\underset{u\in [0,1]}{\sup}~ \underset{x\in F(u)}{\sup}\lvert x \rvert=\infty.$ 
\end{remark}

\begin{definition}
Let $F:[0,1]\rightrightarrows \mathbb{R}$ be a set-valued mapping, then a function $f:[0,1]\rightarrow \mathbb{R}$ is known to be a selection of $F$ if $f(u)\in F(u)$ for each $u\in [0,1].$
\end{definition}

\begin{definition}\cite{aumann1965integrals}
A set-valued map, $F:[a,b]\rightrightarrows \mathbb{R}$ is said to be Borel-measurable if the graph, $Gr(F)$ is a Borel subset of $[a,b]\times \mathbb{R}$. Moreover, if there exists an integrable function, $h:[a,b]\rightarrow \mathbb{R}$ such that $\lvert x \rvert \leq h(u)$ for all $x\in F(u)$, then $F$ is said to be integrably bounded.
\end{definition}
\begin{example}
Let $F:[0,1]\rightrightarrows \mathbb{R}$ be a set-valued map defined as,
$F(u)=[-u,u] \text{ for } u\in [0,1].$
Then, the function, $h:[0,1]\rightarrow \mathbb{R}$ defined as $h(u)=u$  satisfies $\lvert x \rvert \leq h(u)$ for each $x\in F(u)$ and hence $F$ is an integrably bounded function.
\end{example}

\begin{remark}
Every bounded set-valued map is integrably bounded but the converse is not true.\\
Because, if $F:[a,b]\rightrightarrows \mathbb{R}$ is a bounded set-valued map, then by Note \ref{nt2.8} we get $\underset{u\in [a,b]}{\sup}\underset{x\in F(u)}{\sup}\lvert x \rvert = \mathcal{M}$ for some finite $\mathcal{M} \in \mathbb{R}$. Define $h:[a,b]\rightarrow \mathbb{R}$ such that 
\[h(u)=\mathcal{M} \text{ for all } u\in [a,b].\]
Then, $h$ being a constant function it is integrable and satisfies $\lvert x \rvert \leq h(u)$ for each $x\in F(u).$ Hence, $F$ is integrably bounded by $h$.\\
Conversely, assume $F:[0,\infty)\rightrightarrows \mathbb{R}$ is a set-valued map such that $F(u)=[1,1+u)$ for each $u\in [0,\infty)$, then 
\[\sup_{u\in [0,\infty)}\sup_{x\in F(u)} \lvert x \rvert= \sup_{u\in [0,\infty)} (1+u)= \infty.\]
Therefore, $F$ is not bounded but it is integrably bounded by $h:[0,\infty)\rightarrow \mathbb{R}$ defined as $h(u)=1+u$.
\end{remark}

\begin{definition}\cite{CK}
A subset $Y$ of a metric space $X$ is said to be analytic, if $Y$ is a continuous image of a Borel subset of $X$.
A set-valued mapping $F:[a,b]\rightrightarrows \mathbb{R}$ is said to be analytic if its graph is an analytic subset of $[a,b]\times \mathbb{R}$.
\end{definition}

\begin{proposition}\cite{aumann1965integrals}\label{Prop1}
If $F:[a,b]\rightrightarrows \mathbb{R}$ is an analytic set-valued mapping, then there exists a Lebesgue measurable selection of $F$.
\end{proposition}

\subsection*{Set-valued Riemann-Liouville fractional integral } Here, we now introduce the Riemann–Liouville fractional integral of a set-valued mapping via integerable selections.
\begin{definition}
Let $F:[a,b]\rightrightarrows \mathbb{R}$ be a set-valued mapping and $\mathfrak{F}$ be the collection of all integrable selections, $f$ of $F$. Riemann-Liouville set-valued fractional integrable of $F$ is defined as
\[_{a}\mathfrak{J}^{\rho}F(u)=\left\{\frac{1}{\Gamma(\rho)}\int_{a}^{u}(u-t)^{\rho-1}f(t)~\mathrm{d}t : f\in \mathfrak{F}\right\},\] 
where $\rho>0,$ is a real number.
\end{definition}

%Next we shall note some properties of the $ _{a}\mathfrak{J}^{\rho} F$ in the form of theorems. The techniques which we have applied while proving some of the theorems are used in \cite{aumann1965integrals}.

\section{Main Results}
In this section, we establish some fundamental properties of the set-valued Riemann–Liouville fractional integral, which play a crucial role in the development of our main results. 
\begin{theorem} \label{convex}
Let $F:[a,b]\rightrightarrows\mathbb{R}$ be a set-valued mapping with the property that $F(t)$ is convex for each $t\in[a,b]$.
Then for every $\rho>0$ and every $u\in[a,b]$, the set ${}_{a}\mathfrak{J}^{\rho}F(u)$ is convex.
\end{theorem}
\begin{proof}
Fix $u\in[a,b]$. If ${}_{a}\mathfrak{J}^{\rho}F(u)$ is empty or a singleton, there is nothing to prove. Otherwise, take $y_{1},y_{2}\in {}_{a}\mathfrak{J}^{\rho}F(u)$. By the definition of ${}_{a}\mathfrak{J}^{\rho}F(u)$, there exist integrable selections $f$ and $g$ of $F$ such that
$$
y_{1}=\frac{1}{\Gamma(\rho)}\int_{a}^{u}(u-t)^{\rho-1}f(t)dt,
\qquad
y_{2}=\frac{1}{\Gamma(\rho)}\int_{a}^{u}(u-t)^{\rho-1}g(t)dt.
$$
Let $\lambda\in[0,1]$ and define $h=\lambda f+(1-\lambda)g$. Since $f(t),g(t)\in F(t)$ for a.e.~$t$ and each $F(t)$ is convex, we have $h(t)\in F(t)$ for a.e.~$t$. Moreover, $h$ is integrable because it is a linear combination of integrable functions. Hence $h$ is an integrable
selection of $F$, and therefore
$$
\frac{1}{\Gamma(\rho)}\int_{a}^{u}(u-t)^{\rho-1}h(t)dt
\in {}_{a}\mathfrak{J}^{\rho}F(u).
$$
By linearity of the Lebesgue integral,
$$
\frac{1}{\Gamma(\rho)}\int_{a}^{u}(u-t)^{\rho-1}h(t)dt
=
\lambda y_{1}+(1-\lambda)y_{2}.
$$
Thus $\lambda y_{1}+(1-\lambda)y_{2}\in {}_{a}\mathfrak{J}^{\rho}F(u)$ for all $\lambda\in[0,1]$, proving that ${}_{a}\mathfrak{J}^{\rho}F(u)$ is convex. This completes the proof.
\end{proof}


\begin{theorem}\label{THnonempty}
Let $F:[a,b]\rightrightarrows\mathbb{R}$ be a set-valued mapping with nonempty values. Assume that $F$ is Borel-measurable and integrably
bounded. Then for every $\rho>0$ and every $u\in[a,b]$, the set ${}_{a}\mathfrak{J}^{\rho}F(u)$ is nonempty.
\end{theorem}
\begin{proof}
Since $F$ is Borel-measurable with nonempty values, its graph is an
analytic subset of $[a,b]\times\mathbb{R}$. Then by the Proposition \ref{Prop1}, there exists a Lebesgue-measurable selection
$f:[a,b]\to\mathbb{R}$ such that $f(t)\in F(t)$ for a.e.~$t\in[a,b]$. Moreover, since $F$ is integrably bounded, every such measurable
selection is integrable, and hence $f\in\mathfrak{F}$. Fix $u\in[a,b]$. Since $\rho>0$, the kernel $(u-t)^{\rho-1}$ belongs to
$L^{1}([a,u])$. Therefore, the integral
$$ \frac{1}{\Gamma(\rho)}\int_{a}^{u}(u-t)^{\rho-1}f(t)dt $$ is well defined and finite. By the definition of the Riemann-Liouville fractional integral of $F$, this quantity belongs to
${}_{a}\mathfrak{J}^{\rho}F(u)$. Hence
${}_{a}\mathfrak{J}^{\rho}F(u)\neq\emptyset$ for each $u\in[a,b]$.
\end{proof}
\begin{theorem}
Consider $F:[a,b]\rightrightarrows \mathbb{R}$ being a set-valued mapping with nonempty
compact values, Borel measurable and integrably bounded on $[a,b]$, and
let $\rho>0$. If $F$ is bounded on $[a,b]$, then the set-valued
Riemann-Liouville integral
\[
u \longmapsto {}_{a}\mathfrak{J}^{\rho}F(u)
\]
is also bounded on $[a,b]$, that is,
\[
\sup_{u\in[a,b]} H_d\!\left({}_{a}\mathfrak{J}^{\rho}F(u),\{0\}\right) < \infty.
\]
\end{theorem}
\begin{proof}
By using Theorem \ref{THnonempty}, $\mathfrak{F}$ is non-empty and since $F$ is bounded, therefore from Note \ref{NOTE}, we have 
\[\sup_{u\in [a,b]} \sup_{x\in F(u)} \lvert x \rvert < \infty.\]
This means $\underset{u\in [a,b]}{\sup}\lvert f(u) \rvert < \infty$ for each $f\in \mathfrak{F}$. Assume $\mathcal{M}=\underset{f\in \mathfrak{F}}{\sup}\underset{u\in [a,b]}{\sup}f(u)$. Then, we have
\begin{align*}
  H_d\left(~_{a}\mathfrak{J}^{\rho}F(u), \{0\} \right)=& \sup_{u\in [a,b]} \sup_{y\in~_{a}\mathfrak{J}^{\rho}F(u)} \lvert y \rvert\\
  =& \sup_{u\in [a,b]} \sup_{f\in \mathfrak{F}} \left\lvert \frac{1}{\Gamma(\rho)} \int_{a}^{u}(u-t)^{\rho-1}f(t)~\mathrm{d}t \right\rvert\\
  \leq &\frac{1}{\Gamma(\rho)} \sup_{u\in [a,b]}\sup_{f\in \mathfrak{F}} \int_{a}^{u}\lvert (u-t)^{\rho-1} \rvert \lvert f(t) \rvert ~\mathrm{d}t\\
  \leq & \frac{\mathcal{M}}{\Gamma(\rho)}\int_{a}^{u} \lvert (u-t)^{\rho-1}\rvert ~\mathrm{d}t\\
  \leq & \frac{\mathcal{M}}{\Gamma(\rho)} \frac{(u-a)^{\rho}}{\rho}
\end{align*}
Hence,
\[H_d\left(~_{a}\mathfrak{J}^{\rho}F(u), \{0\}\right) \leq \frac{\mathcal{M}}{\Gamma(\rho)}\frac{(b-a)^{\rho}}{\rho}, \text{ for all } u\in [a,b].\]
This completes the proof.
\end{proof}



\begin{theorem}\label{conti}
   Let $F:[a,b]\rightrightarrows\mathbb{R}$ be a set-valued mapping with nonempty compact values, Borel measurable and integrably bounded on $[a,b]$.
Let $\rho>0$. Then, for each $u\in[a,b]$, the set ${}_{a}\mathfrak{J}^{\rho}F(u)$ is nonempty and compact,  and the mapping
$$u \longmapsto {}_{a}\mathfrak{J}^{\rho}F(u)$$
 is continuous on $[a,b]$ with respect to $H_d$.
\end{theorem}
\begin{proof}
Since $F$ is Borel-measurable, integrably bounded and has nonempty compact values, Theorem \ref{THnonempty} guarantees that, for each $u\in[a,b]$, the set of
integrable selections of $F$ is nonempty and therefore
${}_{a}\mathfrak{J}^{\rho}F(u)\neq\emptyset$.
Let $\mathfrak{F}$ denote the family of all integrable selections of $F$. By integrable boundedness, there exists $h\in L^{1}([a,b])$ such that
$$\sup\{|x|:x\in F(t)\}\le h(t)
\quad\text{for a.e. }t\in[a,b].$$
In particular, for each $f\in\mathfrak{F}$ we have
\begin{equation}\label{EQ1}
|f(t)|\le h(t)\quad\text{for a.e. }t\in[a,b].
\end{equation}
For $f\in\mathfrak{F}$ and $u\in[a,b]$, write
$$I_f(u):=\frac{1}{\Gamma(\rho)}\int_{a}^{u}(u-t)^{\rho-1}f(t)dt,$$
so that
$${}_{a}\mathfrak{J}^{\rho}F(u)=\{I_f(u):f\in\mathfrak{F}\}.$$
For any $f\in\mathfrak{F}$ and $u\in[a,b]$, by using \eqref{EQ1} we have
\begin{align*}
    |I_f(u)|
\le
\frac{1}{\Gamma(\rho)}\int_{a}^{u}(u-t)^{\rho-1}|f(t)|dt
\le
\frac{1}{\Gamma(\rho)}\int_{a}^{u}(u-t)^{\rho-1}h(t)dt.
\end{align*}

Since $\rho>0$ and
$h\in L^{1}([a,b])$, the right-hand side is finite. 
Thus there exists $M>0$ such that

\begin{align*}
    |I_f(u)|\le M
\quad\text{for all }u\in[a,b]\text{ and all }f\in\mathfrak{F},
\end{align*}
and therefore ${}_{a}\mathfrak{J}^{\rho}F(u)$ is bounded for each $u$. To see that ${}_{a}\mathfrak{J}^{\rho}F(u)$ is closed, let  $\{I_{f_n}(u)\}_{n\in\mathbb{N}}\subset {}_{a}\mathfrak{J}^{\rho}F(u)$  be a sequence converging to some $y\in\mathbb{R}$.   Since each $f_n$ is an integrable selection of $F$ and $F$ is integrably
bounded, there exists a function $h\in L^{1}([a,b])$ such that
\begin{equation}\label{EQ2}
|f_n(t)|\le h(t)
\qquad\text{for a.e. }t\in[a,b]\text{ and for all }n\in\mathbb{N}.
\end{equation}
Consequently, by passing to a subsequence if necessary, we may assume that
\[
f_n(t)\longrightarrow f(t)
\quad\text{for a.e. }t\in[a,b],
\]
where $f\in L^{1}([a,b])$. Since $f_n(t)\in F(t)$ for a.e.~$t$ and each
$F(t)$ is compact,hence closed, it follows that $f(t)\in F(t)$ for a.e.~$t\in[a,b]$. Therefore, $f$ is also a selection of $F$, that is,
$f\in\mathfrak{F}$.

Now, by \eqref{EQ2} and the fact that the kernel 
$(u-t)^{\rho-1}/\Gamma(\rho)$ is nonnegative and belongs to $L^{1}([a,u])$,
the Dominated Convergence Theorem yields
\[
I_{f_n}(u)
=
\frac{1}{\Gamma(\rho)}\int_{a}^{u}(u-t)^{\rho-1}f_n(t)dt
\;\longrightarrow\;
\frac{1}{\Gamma(\rho)}\int_{a}^{u}(u-t)^{\rho-1}f(t)dt
=
I_f(u).
\]
Since $I_{f_n}(u)\to y$, it follows that $y=I_f(u)$, and hence  $y\in{}_{a}\mathfrak{J}^{\rho}F(u)$.  
Therefore ${}_{a}\mathfrak{J}^{\rho}F(u)$ is closed.  
Being closed and bounded in $\mathbb{R}$, it is compact.  
Together with nonemptiness from Theorem \ref{THnonempty}, this shows that  ${}_{a}\mathfrak{J}^{\rho}F(u)$ is a nonempty compact set for every  $u\in[a,b]$.
Let $u,v\in[a,b]$ and $f\in\mathfrak{F}$ be arbitrary. Without loss of generality assume $u\le v$. Then
\[
\begin{aligned}
I_f(v)-I_f(u)
&=
\frac{1}{\Gamma(\rho)}
\Bigg(
\int_{a}^{v}(v-t)^{\rho-1}f(t)dt
-
\int_{a}^{u}(u-t)^{\rho-1}f(t)dt
\Bigg) \\
&=
\frac{1}{\Gamma(\rho)}
\Bigg(
\int_{a}^{u}\big[(v-t)^{\rho-1}-(u-t)^{\rho-1}\big]f(t)dt
+
\int_{u}^{v}(v-t)^{\rho-1}f(t)dt
\Bigg).
\end{aligned}
\]
From \eqref{EQ1}, we have
$$\big|I_f(v)-I_f(u)\big| \le \Phi(u,v),$$
where
\[
\Phi(u,v) := \frac{1}{\Gamma(\rho)}
\Bigg( \int_{a}^{u}\big|(v-t)^{\rho-1}-(u-t)^{\rho-1}\big|h(t)dt + \int_{u}^{v}(v-t)^{\rho-1}h(t)dt \Bigg).
\]
We now show that $\Phi(u,v)\to 0$ as $v\to u$. For the first integral,
note that for fixed $u\in[a,b]$ and for a.e.~$t\in[a,u)$,
\[ (v-t)^{\rho-1}\to (u-t)^{\rho-1}
\quad\text{as }v\to u, \]
and hence
\[ \big|(v-t)^{\rho-1}-(u-t)^{\rho-1}\big|
\to 0 \quad\text{for a.e. }t\in[a,u]. \] 
 Moreover, since $\rho>0$, we have
\[ \big|(v-t)^{\rho-1}-(u-t)^{\rho-1}\big|
\le (v-t)^{\rho-1}+(u-t)^{\rho-1}, \]
and both $(v-t)^{\rho-1}$ and $(u-t)^{\rho-1}$ belong to $L^{1}([a,u])$ as functions of $t$. Thus
\[ \big|(v-t)^{\rho-1}-(u-t)^{\rho-1}\big|h(t) \le\big[(v-t)^{\rho-1}+(u-t)^{\rho-1}\big]h(t)\in L^{1}([a,u]), \]
and the Dominated Convergence Theorem yields
\[ \int_{a}^{u}\big|(v-t)^{\rho-1}-(u-t)^{\rho-1}\big|\,h(t)dt \longrightarrow 0 \quad\text{as }v\to u. \]

For the second integral, observe that for $v>u$,
\[ 0\le\int_{u}^{v}(v-t)^{\rho-1}h(t)dt \le\int_{u}^{v}(v-t)^{\rho-1}h(t)dt, \] and the right-hand side tends to $0$ as $v\to u$ by the absolute continuity of the Lebesgue integral, since $h\in L^{1}([a,b])$ and the interval $(u,v)$ shrinks to a point. Consequently,
\[ \Phi(u,v)\to 0 \quad\text{whenever }v\to u. \]
Now, fix $u,v\in[a,b]$ and take any $y\in {}_{a}\mathfrak{J}^{\rho}F(u)$. Then
$y=I_f(u)$ for some $f\in\mathfrak{F}$. Let $z:=I_f(v)\in
{}_{a}\mathfrak{J}^{\rho}F(v)$. By the above,
\[ |y-z|=\big|I_f(u)-I_f(v)\big|\le\Phi(u,v), \]
and hence
\[\sup_{y\in {}_{a}\mathfrak{J}^{\rho}F(u)} \inf_{z\in {}_{a}\mathfrak{J}^{\rho}F(v)}|y-z| \le\Phi(u,v). \]
Exchanging the roles of $u$ and $v$ yields
\[ \sup_{z\in {}_{a}\mathfrak{J}^{\rho}F(v)} \inf_{y\in {}_{a}\mathfrak{J}^{\rho}F(u)}|z-y| \le\Phi(u,v). \]
By the definition of the Hausdorff distance, we therefore have
\[ H_d\big({}_{a}\mathfrak{J}^{\rho}F(u),{}_{a}\mathfrak{J}^{\rho}F(v)\big) \le\Phi(u,v)
\quad\text{for all }u,v\in[a,b]. \]
Since $\Phi(u,v)\to 0$ as $v\to u$, this shows that the mapping
\[u\longmapsto {}_{a}\mathfrak{J}^{\rho}F(u) \]
is continuous on $[a,b]$ with respect to the Hausdorff distance $H_d$. This completes the proof. 
\end{proof}
\begin{theorem}\label{BV}
Let $F:[a,b]\rightrightarrows\mathbb{R}$ be a set-valued mapping with nonempty compact convex values, Borel measurable and integrably bounded on $[a,b]$.
Assume that $F$ is of bounded variation on $[a,b]$ with respect to $H_d$ and that $\rho>1$. Then, for each $u\in[a,b]$, the set ${}_{a}\mathfrak{J}^{\rho}F(u)$ is nonempty, compact and convex, and the mapping
$$ u \longmapsto {}_{a}\mathfrak{J}^{\rho}F(u) $$ is of bounded variation on $[a,b]$ with respect to the $H_d$.
\end{theorem}

\begin{proof}
Since $F$ is Borel-measurable, integrably bounded and has nonempty compact
values, Theorem \ref{conti} (applied with $\rho>1$) yields that, for each $u\in[a,b]$, the set ${}_{a}\mathfrak{J}^{\rho}F(u)$ is nonempty and compact and that the mapping $u\mapsto {}_{a}\mathfrak{J}^{\rho}F(u)$ is continuous on $[a,b]$ with respect to $H_d$. Moreover, by Theorem \ref{convex}, since $F$ has convex values, each value ${}_{a}\mathfrak{J}^{\rho}F(u)$ is convex. Thus it remains to show that $u\mapsto {}_{a}\mathfrak{J}^{\rho}F(u)$ has bounded variation on $[a,b]$ with respect to $H_d$.
For each $u\in[a,b]$, the set $F(u)$ is a nonempty compact convex subset of $\mathbb{R}$, hence a closed interval. Define
\[\underline f(u):=\inf F(u),\qquad \overline f(u):=\sup F(u), \]
so that
\[ F(u)=[\underline f(u),\overline f(u)], \qquad u\in[a,b].
\]
Let $\mathcal{P}=\{a=u_0<u_1<\dots<u_n=b\}$ be a partition of $[a,b]$. Then for every $i=1,\dots,n$,
\begin{align*}
    H_d\big(F(u_i),F(u_{i-1})\big) & =
H_d\big([\underline f(u_i),\overline f(u_i)],
        [\underline f(u_{i-1}),\overline f(u_{i-1})]\big) \\&  = \max\Big\{|\underline f(u_i)-\underline f(u_{i-1})|,
          |\overline f(u_i)-\overline f(u_{i-1})|\Big\}.
\end{align*}
Consequently,
\[ \sum_{i=1}^{n}|\underline f(u_i)-\underline f(u_{i-1})|
\le \sum_{i=1}^{n}H_d\big(F(u_i),F(u_{i-1})\big), \]
and similarly for $\overline f$. Taking the supremum over all partitions
$\mathcal{P}$, we obtain
\[ V(\underline f,[a,b])\le V(F,[a,b]),\qquad V(\overline f,[a,b])\le V(F,[a,b])<\infty. \]
In particular, $\underline f$ and $\overline f$ are functions of bounded variation on $[a,b]$ and therefore are bounded on $[a,b]$. For $u\in[a,b]$ define
\[ \big({}_{a}\mathfrak{J}^{\rho}\underline f\big)(u)
:= \frac{1}{\Gamma(\rho)}\int_{a}^{u}(u-t)^{\rho-1}\underline f(t)dt, \qquad
\big({}_{a}\mathfrak{J}^{\rho}\overline f\big)(u)
:= \frac{1}{\Gamma(\rho)}\int_{a}^{u}(u-t)^{\rho-1}\overline f(t)dt. \]
Since $\underline f$ and $\overline f$ are bounded and integrable, the functions ${}_{a}\mathfrak{J}^{\rho}\underline f$ and
${}_{a}\mathfrak{J}^{\rho}\overline f$ are well defined. Now, we claim that both of them are Lipschitz , hence of bounded variation, on $[a,b]$.
Let $g$ be either $\underline f$ or $\overline f$, and set
\[
I_g(u):=\frac{1}{\Gamma(\rho)}\int_{a}^{u}(u-t)^{\rho-1}g(t)dt.
\]
Since $g$ is bounded on $[a,b]$, there exists $M>0$ such that $|g(t)|\le M$
for all $t\in[a,b]$. For $u\in(a,b]$, differentiating under the integral
sign (since $\rho>1$ and the kernel
$(u-t)^{\rho-2}$ is integrable on $[a,u]$) gives
$$I'_g(u)
=
\frac{\rho-1}{\Gamma(\rho)}
\int_{a}^{u}(u-t)^{\rho-2}g(t)dt.$$
Therefore,
\begin{align*}
    |I'_g(u)|
&\le
\frac{\rho-1}{\Gamma(\rho)}
\int_{a}^{u}(u-t)^{\rho-2}|g(t)|dt
\\& \le
\frac{\rho-1}{\Gamma(\rho)}M\int_{a}^{u}(u-t)^{\rho-2}dt \\& =\frac{M}{\Gamma(\rho)}(u-a)^{\rho-1}
\\& \le
\frac{M}{\Gamma(\rho)}(b-a)^{\rho-1}.
\end{align*}
Hence $I'_g$ is bounded on $(a,b]$ and extends continuously to $[a,b]$,
so $I_g$ is Lipschitz on $[a,b]$ with Lipschitz constant
\[
L_0:=\frac{M}{\Gamma(\rho)}(b-a)^{\rho-1}.
\]
Thus both ${}_{a}\mathfrak{J}^{\rho}\underline f$ and
${}_{a}\mathfrak{J}^{\rho}\overline f$ are Lipschitz (in fact with the same
constant $L_0$), and therefore they are functions of bounded variation on
$[a,b]$.

We now show that, for each $u\in[a,b]$,
\begin{equation}\label{EQ3}
{}_{a}\mathfrak{J}^{\rho}F(u)
=
\Big[
\big({}_{a}\mathfrak{J}^{\rho}\underline f\big)(u),
\big({}_{a}\mathfrak{J}^{\rho}\overline f\big)(u)
\Big].
\end{equation}
Let $f\in\mathfrak{F}$ be an integrable selection of $F$, so that
$\underline f(t)\le f(t)\le\overline f(t)$ for all $t\in[a,b]$. Since the
kernel $(u-t)^{\rho-1}/\Gamma(\rho)$ is non negative for $t\in[a,u]$, we obtain
\[
\big({}_{a}\mathfrak{J}^{\rho}\underline f\big)(u)
\le I_f(u)\le
\big({}_{a}\mathfrak{J}^{\rho}\overline f\big)(u),
\] 
which shows that every element of ${}_{a}\mathfrak{J}^{\rho}F(u)$ belongs
to the interval on the right-hand side of \eqref{EQ3}. On the other hand, $\underline f$ and $\overline f$ themselves are selections of
$F$, so the corresponding integrals
${}_{a}\mathfrak{J}^{\rho}\underline f(u)$ and
${}_{a}\mathfrak{J}^{\rho}\overline f(u)$ belong to
${}_{a}\mathfrak{J}^{\rho}F(u)$. Hence the interval in
\eqref{EQ3} is contained in ${}_{a}\mathfrak{J}^{\rho}F(u)$,
and equality follows.
Thus, for every $u\in[a,b]$,
\begin{equation*}
    {}_{a}\mathfrak{J}^{\rho}F(u)
=
\big[ A(u),B(u)\big],
\qquad \text{where}\quad
A(u):={}_{a}\mathfrak{J}^{\rho}\underline f(u),\quad
B(u):={}_{a}\mathfrak{J}^{\rho}\overline f(u).
\end{equation*}
In particular, ${}_{a}\mathfrak{J}^{\rho}F(u)$ is compact and convex. Let $\mathcal{P}=\{a=u_0<\dots<u_n=b\}$ be a partition of $[a,b]$. Then for each $i=1,\dots,n$,
\begin{align*}
    H_d\big({}_{a}\mathfrak{J}^{\rho}F(u_i),{}_{a}\mathfrak{J}^{\rho}F(u_{i-1})\big)
&= H_d\big([A(u_i),B(u_i)],[A(u_{i-1}),B(u_{i-1})]\big) \\& =\max\big\{|A(u_i)-A(u_{i-1})|,
         |B(u_i)-B(u_{i-1})|\big\}.
\end{align*}
Hence
$$\sum_{i=1}^{n}
H_d\big({}_{a}\mathfrak{J}^{\rho}F(u_i),{}_{a}\mathfrak{J}^{\rho}F(u_{i-1})\big)
\le
\sum_{i=1}^{n}\Big(|A(u_i)-A(u_{i-1})|
                   +|B(u_i)-B(u_{i-1})|\Big).$$
Taking the supremum over all partitions $\mathcal{P}$, we obtain
\[
V\big({}_{a}\mathfrak{J}^{\rho}F,[a,b]\big)
\le
V(A,[a,b])+V(B,[a,b])<\infty,
\]
since $A$ and $B$ are of bounded variation. Thus
$u\mapsto{}_{a}\mathfrak{J}^{\rho}F(u)$ is of bounded variation on $[a,b]$ with respect to $H_d$. This completes the proof.
\end{proof}

\begin{theorem} \label{LP}
Let $F:[a,b]\rightrightarrows\mathbb{R}$ be a set-valued mapping with nonempty compact values, Borel measurable and integrably bounded on $[a,b]$.
Assume that $F$ is Lipschitz on $[a,b]$ with respect to $H_d$ and that $\rho>1$. Then, for each $u\in[a,b]$, the set ${}_{a}\mathfrak{J}^{\rho}F(u)$ is nonempty and compact, and the mapping
\[
u \longmapsto {}_{a}\mathfrak{J}^{\rho}F(u)
\]
is Lipschitz on $[a,b]$ with respect to $H_d$.
\end{theorem}
    \begin{proof}
By Theorem \ref{conti} (with $\rho>1$), for each $u\in[a,b]$ the set ${}_{a}\mathfrak{J}^{\rho}F(u)$ is nonempty and compact. Thus, it remains to show that the mapping $u\mapsto{}_{a}\mathfrak{J}^{\rho}F(u)$ is Lipschitz with respect to $H_d$.     
Let $L_F\ge0$ be a Lipschitz constant of $F$ with respect to $H_d$, that is
$$H_d\big(F(u),F(v)\big)\le L_F|u-v|
\quad\text{for all }u,v\in[a,b].$$


To prove the Lipschitz continuity of ${}_{a}\mathfrak{J}^{\rho}F$, we first  establish a uniform bound on the values of $F$. Fix $u_{0}\in[a,b]$ and  choose any $x_{0}\in F(u_{0})$. Since $F(u_{0})$ is compact, it is bounded,  and therefore 
\[
R_{0}:=\sup\{\,|x|:\,x\in F(u_{0})\,\}<\infty.
\]
For arbitrary $u\in[a,b]$ and $x\in F(u)$,by the definition
of the Hausdorff distance there exists $y\in F(u_0)$ such that
$$|x-y|\le H_d\big(F(u),F(u_0)\big).$$
Because the Lipschitz continuity of $F$, we have
$$|x-y|\le H_d\big(F(u),F(u_0)\big)\le L_F|u-u_0|\le L_F(b-a).$$
Since $y\in F(u_{0})$, we have $|y|\le R_{0}$, and hence
\[
|x|\le |y|+|x-y|\le R_{0}+L_{F}(b-a)=:M.
\]
Thus,
\[
\sup\{\,|x|:\,x\in F(u)\,\}\le M
\qquad\text{for all }u\in[a,b].
\]
Consequently, every integrable selection $f$ of $F$ satisfies
\begin{equation}\label{EQ4}
|f(t)|\le M
\qquad\text{for all }t\in[a,b].
\end{equation}


For such a selection $f$, define
$$
I_f(u):=\frac{1}{\Gamma(\rho)}\int_{a}^{u}(u-t)^{\rho-1}f(t)dt,
\qquad u\in[a,b].
$$
As in the proof of the above theorem,  since $\rho>1$ and $f$ is bounded, we may differentiate under the integral sign to obtain, for
$u\in(a,b]$,
$$
I_f'(u)
=
\frac{\rho-1}{\Gamma(\rho)}
\int_{a}^{u}(u-t)^{\rho-2}f(t)dt.
$$
Using \eqref{EQ4}
$$|I_f'(u)|
\le
\frac{\rho-1}{\Gamma(\rho)}M\int_{a}^{u}(u-t)^{\rho-2}dt
=
\frac{M}{\Gamma(\rho)}(u-a)^{\rho-1}
\le
\frac{M}{\Gamma(\rho)}(b-a)^{\rho-1}
=:L_0.$$
Hence $I_f$ is Lipschitz on $[a,b]$ with Lipschitz constant at most $L_0$, and the constant $L_0$ is independent of the particular selection $f$.
Now let $u,v\in[a,b]$ and $f\in\mathfrak{F}$ be an integrable selection of
$F$. Then
\[
|I_f(u)-I_f(v)|\le L_0|u-v|.
\]
If $y\in {}_{a}\mathfrak{J}^{\rho}F(u)$, then $y=I_f(u)$ for some
$f\in\mathfrak{F}$, and taking $z:=I_f(v)\in{}_{a}\mathfrak{J}^{\rho}F(v)$
we obtain
\[
|y-z|\le L_0|u-v|.
\]
Therefore
\[
\sup_{y\in {}_{a}\mathfrak{J}^{\rho}F(u)}
\inf_{z\in {}_{a}\mathfrak{J}^{\rho}F(v)}|y-z|
\le L_0|u-v|.
\]
By symmetry (interchanging $u$ and $v$), we also have
\[
\sup_{z\in {}_{a}\mathfrak{J}^{\rho}F(v)}
\inf_{y\in {}_{a}\mathfrak{J}^{\rho}F(u)}|z-y|
\le L_0|u-v|.
\]
Hence, by the definition of the Hausdorff distance,
\[
H_d\big({}_{a}\mathfrak{J}^{\rho}F(u),
        {}_{a}\mathfrak{J}^{\rho}F(v)\big)
\le L_0|u-v|
\quad\text{for all }u,v\in[a,b],
\]
which shows that $u\mapsto{}_{a}\mathfrak{J}^{\rho}F(u)$ is Lipschitz on $[a,b]$ with respect to $H_d$. This completes the proof.
\subsection*{Exitence of regular selections:}
\end{proof}
\begin{theorem} \label{Maintheorem}
    Let $F:[a,b]\rightrightarrows\mathbb{R}$ be a set valued mapping with nonempty
compact values, Borel measurable and integrably bounded on $[a,b]$, and let
$\rho>1$. Let ${}_{a}\mathfrak{J}^{\rho}F$ be the Riemann-Liouville set valued fractional integral of $F$. Then the following hold.
    \begin{itemize}
        \item[(a)] If $F$ has compact convex values and is of bounded variation on $[a,b]$ with respect to $H_d$, then the mapping
$$ u \longmapsto {}_{a}\mathfrak{J}^{\rho}F(u)$$ admits a continuous selection which is itself of bounded variation.
        \item[(b)] If $F$ is Lipschitz on $[a,b]$ with respect to $H_d$, then the mapping
$$ u \longmapsto {}_{a}\mathfrak{J}^{\rho}F(u)$$ admits a Lipschitz selection.
    \end{itemize}
\end{theorem}
\begin{proof} Now, the Theorems \ref{THnonempty}, \ref{conti}, \ref{BV}, \ref{LP} and Theorem 3 of
Belov and Chistyakov \cite{SAVV} will be used in the proof of this theorem.\\
Let $\rho>1$ and let $F:[a,b]\rightrightarrows\mathbb{R}$ satisfy the assumptions of Theorem \ref{Maintheorem}.  Define the set-valued mapping
$$G(u):={}_{a}\mathfrak{J}^{\rho}F(u), \qquad u\in[a,b].$$ By Theorem \ref{THnonempty} for each $u\in[a,b]$ the set $G(u)$ is nonempty and  compact. 
\begin{itemize}
    \item[(a)] From Theorems \ref{conti} and \ref{BV}, $G$ is continuous and of bounded variation on $[a,b]$ with respect to Hausdorff distance $H_d$. Theorem 3 of Belov and Chistyakov \cite{SAVV} guarantees the existence of a regular selection $g:[a,b]\to\mathbb{R}$ of $G$, that is,
$$g(u)\in G(u)={}_{a}\mathfrak{J}^{\rho}F(u)
\quad\text{for all }u\in[a,b],$$
and $g$ is continuous and of bounded variation on $[a,b]$ with 
$$V(g,[a,b]) \le V(G,[a,b]).$$

    \item[(b)] From Theorem \ref{LP}, $G$ is Lipschitz on $[a,b]$ with respect to the Hausdorff metric $H_g$.  Theorem 3 of Belov and Chistyakov \cite{SAVV} yields a 
Lipschitz selection $g:[a,b]\to\mathbb{R}$ of $G$, i.e.,
$$g(u)\in G(u)={}_{a}\mathfrak{J}^{\rho}F(u)
\quad\text{for all }u\in[a,b],$$
 with
$$\operatorname{Lip}(g)\ \le\ \operatorname{Lip}(G).$$
\end{itemize}
This completes the proof.
\end{proof}
\begin{remark}
The convexity of the values of $F$ is only used in the proof of the bounded variation property of the Riemann–Liouville fractional integral in Theorem \ref{BV}. The subsequent existence of regular selections via the result of Belov and Chistyakov \cite{SAVV} does not require convexity.
\end{remark}
\begin{theorem}[Extremal selections of ${}_{a}\mathfrak{J}^{\rho}F$]

Let $F:[a,b]\rightrightarrows\mathbb{R}$ be a set-valued mapping with nonempty compact values, and let $\rho>0$. For each $u\in[a,b]$, let
$$G(u):={}_{a}\mathfrak{J}^{\rho}F(u).$$
By Theorem \ref{THnonempty}, $G(u)$ is a nonempty compact subset of $\mathbb{R}$ for each $u\in[a,b]$. Define the extremal selections
$$ g_{-}(u):=\min G(u),  \qquad g_{+}(u):=\max G(u), \qquad u\in[a,b]. $$
Then the following assertions hold:
\begin{itemize}
    \item[(a)] Under the assumptions of Theorem \ref{conti},  $G$ is continuous on $[a,b]$, and both $g_{-}$ and $g_{+}$ are continuous  selections of $G$ on $[a,b]$.
    
    \item[(b)] Under the assumptions of Theorem \ref{BV} (in particular $\rho>1$) 
$G$ is of bounded variation on $[a,b]$, and both $g_{-}$ and $g_{+}$ are of 
bounded variation on $[a,b]$. Moreover,
$$ V\big(g_{-},[a,b]\big)\le V\big(G,[a,b]\big), \qquad V\big(g_{+},[a,b]\big)\le V\big(G,[a,b]\big). $$
\item[(c)] Under the assumptions of Theorem \ref{LP} (in particular $\rho>1$), $G$ is 
Lipschitz on $[a,b]$, and both $g_{-}$ and $g_{+}$ are Lipschitz on $[a,b]$. 
In addition,
$$ \operatorname{Lip}(g_{-})\le \operatorname{Lip}(G), \qquad \operatorname{Lip}(g_{+})\le\operatorname{Lip}(G).
$$
\end{itemize}
\end{theorem}
    \begin{proof}
By Theorem \ref{THnonempty}, each $G(u)$ is a nonempty compact subset of $\mathbb{R}$, hence $\min G(u)$ and $\max G(u)$ are well defined and yield selections $g_{-}$ and $g_{+}$. The continuity, bounded variation, and Lipschitz regularity of $g_{-}$ and $g_{+}$ follow directly from the corresponding regularity of $G$ (from Theorems \ref{conti}, \ref{BV} and \ref{LP}) and standard properties of the Hausdorff distance on $\mathbb{R}$.
\end{proof}



\section*{\textbf{Application}}
One of the important applications of research presented in this paper is in the area of differential inclusion. In order to prove existence and uniqueness, one can directly utilize the selection theorem.   \\
 Let $\alpha\in(1,2)$ and consider the Caputo fractional
differential inclusion
\begin{equation}\label{CPEQN}
{}^{C}D^{\alpha}u(t)\in F(t,u(t))
\quad\text{for a.e. }t\in[t_{0},T],
\qquad
u(t_{0})=u_{0},\quad u'(t_{0})=u_{1},
\end{equation}
where ${}^{C}D^{\alpha}$ denotes the Caputo derivative of order $\alpha$ and
$F:[t_{0},T]\times\mathbb{R}\rightrightarrows\mathbb{R}$ is a set-valued
mapping with nonempty compact values.

 The existence of a solution  usually requires that $F(t, u)$ be measurable in $t$, upper hemicontinuous function of $u$, and $F(t, u)$ a closed, convex set for all $(t,u)$. The uniqueness requires the Lipschitz condition on $F$, i.e. Hausdorff distance
 $$H_d(F(t, u_1), F(t, u_2)) \leq L \|u_1 - u_2\|$$ for some $L > 0$. For more details, we refer to \cite{incl}. The existence of a solution is equivalent to looking for fixed points of the corresponding integral form. We look for a solution in the set $$_{a}\mathfrak{J}^{\rho}F(t)=\left\{\frac{1}{\Gamma(\rho)}\int_{a}^{t}(t-s)^{\rho-1}f(s)~\mathrm{d}s : f\in \mathfrak{F}\right\}.$$ 
 It is well known that problem \eqref{CPEQN} is equivalent to
the integral inclusion
\begin{equation}\label{INGEQN}
u(t)\in u_{0}+u_{1}(t-t_{0})
+\frac{1}{\Gamma(\alpha)}
\int_{t_{0}}^{t}(t-s)^{\alpha-1}v(s)\,ds,
\qquad v(s)\in F(s,u(s))\ \text{a.e.}
\end{equation}
for $t\in[t_{0},T]$.

Setting $\rho=\alpha>1$, the integral term in \eqref{INGEQN}
coincides with the Riemann-Liouville fractional integral of order $\rho$
applied to a selection of the mapping $s\mapsto F(s,u(s))$. By using the results proved, it is evident that a selection exists and hence a solution exists. Hence, we can directly apply the results on the problem of differential inclusions of fractional order to prove the existence and uniqueness of the solution.
\\
\section{Future remarks}
The computation of fractal dimensions has long been a central and fascinating topic in fractal theory. In the past few years, the study of fractal dimension (such as box dimension and the Hausdorff dimension) of the graph of the integral of a function is getting the attention of the researchers. For example, Liang \cite{L1} introduced the box dimension of the graph of the Riemann-Liouville fractional integral of a continuous function of bounded variation on a closed bounded interval. Subsequently, several significant developments in this direction have been reported; see, for instance, \cite{SS, SSS, SS1, L2, MV, SV,SV1, BY, BY1}. Motivated by these advances, it would be natural and interesting to investigate analogous dimensional properties for set-valued Riemann–Liouville fractional integrals, including the fractal dimensions of their graphs under suitable regularity assumptions on the underlying set-valued mappings.
\subsection*{Acknowledgements} The idea of this work originated during the workshop on Fractals: GAANA 2025 organized by IIIT Prayagraj.




  
\bibliographystyle{amsplain}
\begin{thebibliography}{10}
\bibitem{aumann1965integrals} R.J. Aumann, Integrals of set-valued functions, J. Math. Anal. Appl. 12 (1) (1965) 1–12.
\bibitem{art} Z. Artstein, Set-Valued Measures. Transactions of the American Mathematical Society, 165,  (1972) 103–125.
\bibitem{incl} J. P. Aubin, A.  Cellina,  Differential inclusions: Set-Valued Maps and Viability Theory, Grundlehren Math. Wiss., 264, Springer-Verlag, Berlin, (1984).

\bibitem{aubin1} J. P. Aubin,  H.  Frankowska, Set-Valued Analysis. Birkhauser (1990), ISBN
978-0-8176-4847-3.

\bibitem{abbas2} S. Abbas, Existence of solutions to fractional order ordinary and delay differential equations and applications, Electron. J. Differential Equations (09) (2011) pp. 11.

\bibitem{lup1} U. Abbas, V. Lupulescu, D. O’Regan, Y Awais,  Neutral set differential equations. Czech Math J 65, 593–615 (2015). 
\bibitem{jehad1}  J. Alzabut, M. Houas,  M. I. Abbas, Application of fractional quantum calculus on coupled hybrid differential systems within the sequential Caputo fractional  q-derivatives, Demonstr. Math. 56 (2023).

\bibitem{azam} A. Abbas, A. Azam, N. Mehmood, F. Ali, Existence of Solutions to Fractional Differential Inclusions with Non-Separated and Multipoint Boundary Constraints. Fractal and Fractional. 2025; 9(9):577.

\bibitem{SAVV}  S. A. Belov, V.V Chistyakov, A selection principle for mappings of bounded variation. J. Math. Anal. Appl. 249, 351-366 (2000)

\bibitem{GB} G. Beer, Metric spaces on which continuous functions are uniformly continuous and Hausdorff distance, Proc. Amer. Math. Soc. 95 (1985) 653-658.
\bibitem{anca} A. Croitoru, R. Mesiar, A. R. Sambucini, B. Satco, Special Issue on Set Valued Analysis (2021).
\bibitem{hess} H. Christian, CHAPTER 14 - Set-Valued Integration and Set-Valued Probability Theory: An Overview, Handbook of Measure Theory, Volume I, 2002, Pages 617, 619-673.
\bibitem{SS} S. Chandra and S. Abbas, Box dimension of mixed Katugampola fractional integral of two-dimensional
continuous functions, Fract. Calc. Appl. Anal. 25 (2022) 1022–1036.
\bibitem{SSS} S. Chandra, S. Abbas, Analysis of fractal dimension of mixed Riemann-Liouville integral, Numer. Algorithms 91 (2022) 1021–1046.
\bibitem{SS1} S. Chandra, S Abbas, Y. S. Liang, A note on fractal dimension of Riemann–Liouville fractional integral. Fractals 2024;32:2440001.
\bibitem{huang} Y. Huang, J. Lv, Z. Liu,  Existence results for Riemann-Liouville fractional evolution inclusions, Miskolc Math. Notes 17 (2016), no. 1, 305–325.
\bibitem{frac2} H. Khan, S. Ahmed, J. Alzabut, A. T. Azar, A generalized coupled system of fractional differential equations with application to finite time sliding mode control for Leukemia therapy, Chaos Solitons Fractals 174 (2023), Paper No. 113901, 12 pp.

\bibitem{CK} C. Kuratowski, \emph{Topologie I}, Monografie Matematyczne, Vol.~20,
Polish Mathematical Society, Warsaw, 1948.

\bibitem{kara} H. Kara, M. A. Ali, H. Budak, Hüseyin, Hermite-Hadamard-Mercer type inclusions for interval-valued functions via Riemann-Liouville fractional integrals, Turkish J. Math. 46 (2022), no. 6, 2193–2207.
\bibitem{L1} Y. S. Liang, Box dimensions of Riemann-Liouville fractional integrals of continuous functions of bounded variation, Nonlin. Anal. 72(11) (2010) 4304-4306. 
\bibitem{L2} Y. S. Liang, Fractal dimension of Riemann-Liouville fractional integral of 1-dimensional continuous functions. Fract. Calc. Appl. Anal. 21(6), 1651–1658 (2019)
\bibitem{li2}H. Li, C.  Huang, T. Li, Dynamic complexity of a fractional-order predator-prey system with double delays, Phys. A 526 (2019), 120852, 16 pp.
	
\bibitem{lup} V. Lupulescu, Fractional calculus for interval-valued functions,  Fuzzy Sets and Systems
Volume 265, 15 April 2015, Pages 63-85.
\bibitem{pod} I. Podlubny,  Fractional differential equations. An introduction to fractional derivatives, fractional differential equations, to methods of their solution and some of their applications, Math. Sci. Engrg., 198, Academic Press, Inc., San Diego, CA, (1999).
 \bibitem{ergen} D. M. Riley, A. Stathas, D. Gutiérrez-Oribio, I. Stefanou,  Reinforcement Learning-Based Adaptive Time-Integration for Nonsmooth Dynamics. arXiv:2501.08934 (2025).

\bibitem{sriv} H. M. Srivastava, S. K. Sahoo, P. O. Mohammed, B. Kodamasingh, Y. S. Hamed, New Riemann–Liouville Fractional-Order Inclusions for Convex Functions via Interval-Valued Settings Associated with Pseudo-Order Relations. Fractal and Fractional. 2022; 6(4):212.


\bibitem{shen} Y. Shen, The Cauchy-type problem for interval-valued fractional differential equations with the Riemann-Liouville gH-fractional derivative. Adv Differ Equ 2016, 102 (2016).


\bibitem{stef} L. Stefanini, B. Bede, Generalized Hukuhara differentiability of interval-valued functions and interval differential equations, Nonlinear Anal., Theory Methods Appl., 71 (3–4) (2009), pp. 1311-1328.
\bibitem{MV} M. Verma, A. Priyadarshi and S. Verma, Vectorvalued fractal functions: Fractal dimension and fractional calculus, Indag. Math. 34(4) (2023) 830–853.
\bibitem{SV} S. Verma and P. Viswanathan, A note on Katugampola fractional calculus and fractal dimensions,
Appl. Math. Comput. 339 (2018) 220–230.
\bibitem{SV1} S. Verma, P. Viswanathan, Bivariate functions of bounded variation: Fractal dimension and fractional
integral. Indag. Math. 31(2), 294–309 (2020).
\bibitem{BY} B. Y. Yu, Y. S. Liang, Approximation with continuous functions preserving fractal dimensions of the Riemann–Liouville operators of fractional calculus. Fract Calc Appl Anal 2023;26:2805–36.
\bibitem{BY1} B. Y. Yu, B Selmi, Y. S., Liang, General fractal dimensions of graphs of continuous functions associated with the Katugampola fractal integral. Fractals 2550040 (2025)
\bibitem{frac1} A. Zada, S. Ali,  T. Li, Tongxing, Analysis of a new class of impulsive implicit sequential fractional differential equations, Int. J. Nonlinear Sci. Numer. Simul. 21 (2020), no. 6, 571–587.

\end{thebibliography}
	
 \end{document}

%------------------------------------------------------------------------------
% End of journal.tex
